\documentclass[11pt,a4paper,DIV9,german]{scrartcl}
%%%%%%%%%%%%%%%%%%%%%%%%%%%%%%%%%%%%%%%%%%%%%%%%%%%%%%%%%%%%%%%%%%%%%%%%%%
% Anpassungen f"ur das jeweilige "Ubungsblatt
\def\blnr{5}              % Blattnummer
\def\sdate{29.3.2019,}    % Abgabedatum
\def\exnr{1}              % Nummer der ersten Aufgabe des aktuellen Blattes
%%%%%%%%%%%%%%%%%%%%%%%%%%%%%%%%%%%%%%%%%%%%%%%%%%%%%%%%%%%%%%%%%%%%%%%%%%
\usepackage{babel,theorem}
\usepackage{exscale}
\usepackage{amsmath,amsfonts,amssymb,dsfont}
\usepackage{a4wide}
\usepackage{epic,eepic}
\usepackage{graphicx} % einbinden von Grafiken mit \includegraphics
\usepackage{url}

\pagestyle{empty}
\parindent0em
\parskip1ex
\addtolength{\topmargin}{0cm}
\addtolength{\oddsidemargin}{0cm}
\addtolength{\textheight}{2cm}
\addtolength{\textwidth}{0cm}

% neuer Z"ahler f"ur die Aufgaben
\newcounter{aufgnr}
\setcounter{aufgnr}{\exnr}
\addtocounter{aufgnr}{-1}

% in enumerate-Umgebungen erst Kleinbuchstaben, dann kleine r"omische Zahlen
\renewcommand{\labelenumi}{(\alph{enumi})\,}
\renewcommand{\labelenumii}{(\roman{enumii})\,}

% Stil f"ur Definitionen, S"atze usw.
\theoremstyle{break}   % Titelzeile abgesetzt
\newtheorem{Satz}{Satz}
\newtheorem{Lem}{Lemma}
\newtheorem{Def}{Definition}
\theorembodyfont{\upshape}  % normaler Font f"ur Beispiel, Aufgabe
\newtheorem{Bsp}{Beispiel}
\newtheorem{aufg}[aufgnr]{Aufgabe}
\newenvironment{Aufgabe}[1][z*+]%
{\begin{aufg}}
  {\end{aufg}\vspace{1.5ex}}
\newtheorem{praktaufg}[aufgnr]{Aufgabe^*}
\newenvironment{PraktAufgabe}[1][]%  der optionale Param. muss in [[ .. ]]
{\begin{praktaufg}#1}
  {\end{praktaufg}\vspace{1.5ex}}

\newcommand{\zweifaelles}[3]{ \left\{ \begin{array}{ccc} #1 & \hspace{.8cm}  \mbox{ falls } #2 \, , &     \\[1ex]
                               #3 & \mbox{ sonst } &     \end{array} \right.     }
% Mathematische K"urzel
\def\R{\mathbb{R}}
\def\S{\mathbb{S}}
\def\C{\mathbb{C}}
\def\Q{\mathbb{Q}}
\def\N{\mathbb{N}}
\def\Z{\mathbb{Z}}
\def\P{\mathbb{P}}
\def\K{\mathbb{K}}
\def\O{{\mathcal O} }
\def\G{{\mathcal G} }
\def\T{{\mathcal T} }
\def\U{{\mathcal U} }

\newcommand{\arrl}{\begin{array}{l}}
\newcommand{\are}{\end{array}}
\newcommand{\arrc}{\begin{array}{c}}
\newcommand{\mza}{\begin{eqnarray}}
\newcommand{\mze}{\end{eqnarray}}
\newcommand{\dn}{\frac{\partial}{\partial n}}
\newcommand{\dny}{\frac{\partial}{\partial n(y)}}
\newcommand{\dd}[2]{\ensuremath{\frac{\partial^{#1}}{\partial #2}}}
\newcommand{\ov}{\overline}
\newcommand{\ma}{\mid\!}
\newcommand{\me}{\!\mid}
\newcommand{\impl}{\Rightarrow}
\newcommand{\sca}[3]{\ensuremath{(#1,#2)_{#3}}}
\newcommand{\nor}[2]{\ensuremath{\parallel\! #1 \!\parallel_{#2}}}
\newcommand{\ra}{\rightarrow}
\newcommand{\phin}{\varphi}
\newcommand{\diver}{\mbox{div}}

\def\eps{{\varepsilon}}
\def\i{{\sl i\;}}
\def\e{{\rm e}}
\def\D{{\rm D\;}}
\def\norm{|\!|\!|}

\DeclareMathOperator{\spann}{span}
\DeclareMathOperator{\cond}{cond}
\DeclareMathOperator{\round}{round}
\DeclareMathOperator{\supp}{supp}

\begin{document}

\parbox{0ex}{ \scalebox{.27}{  \includegraphics[width=.5\textwidth,angle=-90]{cm-logo2.eps} }   }   \\

\parbox{25ex}{
  Prof.~Dr.~S.~Sauter\\
  Institut f"ur Mathematik\\
  Universit"at Z"urich
  }
%
\rule[0cm]{0.cm}{.01cm}
\hfill  \parbox{0.6\textwidth}{
  {\sf\LARGE\bfseries Numerik I}
  {\sf\Large\bfseries \;\;--\;\; Homework \blnr }\\[1.5ex]
  Deadline: \sdate\ 13:00
  }



\vspace{10ex}


% \scalebox{.27}{ \put(1250,0){ \includegraphics[width=.5\textwidth,angle=-90]{cm-logo2.eps} }  }


%%%%%%%%%%%%%%%%%%%%%%%%%%%%%%%%%%%%%%%%%%%%%%%%%%%%%%%%%%%%%%%%%%%%%%%%%%%
%
% Aufgaben
%
%%%%%%%%%%%%%%%%%%%%%%%%%%%%%%%%%%%%%%%%%%%%%%%%%%%%%%%%%%%%%%%%%%%%%%%%%%%
\paragraph{Restriction}
The use of built-in functions for Fourier transform is not allowed.\\

\textbf{Exercise 1} (6 P.) (Programming task)\\
Write a function that implements the Fourier transform given by {Definition 5.2} in the lecture notes.  The function should take the real-valued vector \(\left(f_k\right)_{k=0}^{N-1}\) as the input argument and return as the output
\begin{enumerate}
  \item the Fourier coefficients \(\hat{f}_k\);
  \item the coefficients \(\hat{A}_m\) and \(\hat{B}_m\) of the trigonometric polynomial \(T(t)\) for real-valued \(f_k\) (Equation 5.5 in the lecture notes);
  \item a plot of the trigonometric polynomial \(T(t)\).
\end{enumerate}


\textbf{Exercise 2} (6 P.) (Programming task)\\
On the course homepage there is a data file \textit{f.dat} which stores the evaluations of a continuous function \(f:[0,2x_0]\rightarrow\mathds{R}\)
\begin{equation*}
  f_k=f\left(\frac{2x_0k}{N}\right)\quad\mathrm{where}~k=0,1,\cdots,N-1
\end{equation*}
for \(N=2^{16}\).  The function has the property \(f(0)=f(2x_0), \; x_0 = 
101.15520937701599\).

For the refinement level \(l=2,4,6,8,10,12,14,16\),
\begin{enumerate}
  \item Use your program in Exercise 1 to calculate the trigonometric polynomial \(T(t)\) by using the following data sets respectively as the input argument
  \begin{equation*}
    f\left(\frac{2x_0k}{2^l}\right)\quad\mathrm{where}~k=0,1,\cdots,2^l-1;
  \end{equation*}
  and show the trigonometric polynomials in the plot.
 
  \item Calculate the approximation error for each data set
  \begin{equation*}
    \max_{0\leq k \leq N-1}\left|T_l\left(\frac{2x_0k}{N}\right)-f_k\right|
  \end{equation*}
  and record them in a table.
\end{enumerate}

\end{document}